\section{Computational Study}

The comparison is based on a benchmark runned on the same 100 instances for all four formulations:
\emph{undirected cut} (UC), \emph{undirected flow} (UF),
\emph{directed cut} (DC), and \emph{directed flow} (DF).
All runs used a time limit of 120 seconds, one thread, seed 0, and the same
instance set. Therefore the comparison is fair instance-by-instance.

The main result is very clear: \textbf{the directed flow formulation dominates the
other three formulations on this benchmark set}. After correcting the DF
implementation, it solves all 100 instances to optimality and is the fastest
formulation on almost every instance.

\subsection{Overall solution time}

Table~\ref{tab:exercise4-summary} summarizes the main performance indicators.

\begin{table}[h!]
\centering
\small
\setlength{\tabcolsep}{4pt}
\begin{tabular}{lrrrrrrrr}
\toprule
Form. & Opt. & TL & Med.\ time & Mean time & Med.\ gap & Med.\ B\&B & Med.\ sep. & Med.\ sep. \\
& & & (s) & (s) & & nodes & time (s) & share \\
\midrule
DC & 100 & 0  & 4.045  & 14.203 & 0.094 & 48.0   & 0.725  & 0.242 \\
DF & 100 & 0  & 0.292  & 4.617  & 0.003 & 1.0    & 0.000  & 0.000 \\
UC & 22  & 78 & 120.016 & 99.506 & 0.233 & 1816.5 & 18.467 & 0.218 \\
UF & 39  & 61 & 120.096 & 79.769 & 0.147 & 770.0  & 0.000  & 0.000 \\
\bottomrule
\end{tabular}
\caption{Main benchmark summary on 100 instances.}
\label{tab:exercise4-summary}
\end{table}

% Placeholder: add a runtime boxplot here if desired.
% Example file from this project:
% test_output/plots/runtime_boxplot.svg
%
% \begin{figure}[h!]
% \centering
% \includegraphics[width=0.8\textwidth]{<path-to-runtime-plot>}
% \caption{Runtime distribution by formulation.}
% \label{fig:runtime-boxplot}
% \end{figure}

The most important conclusions are:

\begin{itemize}
  \item DF and DC are the only formulations that solve \textbf{all 100 instances}.
  \item DF has by far the best median runtime: \SI{0.292}{s}. DC is second with
  \SI{4.045}{s}, while UF and UC still have median runtimes essentially equal
  to the time limit, which means they time out on at least half of the instances.
  \item In pairwise comparisons, DF beats DC on 99 instances, UF on 99 instances,
  and UC on all 100 instances.
  \item If we look at the fastest formulation per instance, DF wins
  \textbf{98 out of 100} instances. DC wins 1, UF wins 1, and UC wins none.
\end{itemize}

So with respect to overall runtime, the ranking is:
\[
\text{DF} \gg \text{DC} \gg \text{UF} > \text{UC}.
\]

\subsection{Quality of LP at root node}

The quality of the root relaxation is measured here by the root LP gap.
The smaller the better, because it means the root bound is closer to the final
optimal objective value.

The median root LP gaps are:

\begin{itemize}
  \item DF: 0.0031
  \item DC: 0.0938
  \item UF: 0.1472
  \item UC: 0.2335
\end{itemize}

Taken at face value, this ranking matches the runtime ranking almost perfectly.
DF has the strongest reported root relaxation, and UC has by far the weakest.
This is important because a
stronger root bound usually reduces the amount of branch-and-bound search that
is needed later, and that is exactly what we observe in the benchmark.

However, this part should be interpreted with some care. For the corrected DF
model, many instances are solved so quickly that the root-node callback does
not always record a root bound. Therefore the reported DF root-gap values are
available only for a subset of instances. So the value 0.0031 is still
informative, but it is not as directly comparable to the other three
formulations as the runtime and branch-and-bound statistics are.

Even with that caveat, the overall picture is still very favorable to DF. Its
reported root gaps are extremely small whenever they are recorded, while DC is
clearly second-best and UC remains the weakest formulation at the root.

\begin{figure}[h!]
\centering
\includegraphics[width=0.8\textwidth]{imgs_final/root_gap_boxplot.png}
\caption{Root LP gap distribution by formulation.}
\end{figure}

\subsection{Number of branch-and-bound nodes}

The median number of branch-and-bound nodes is:

\begin{itemize}
  \item DF: 1.0
  \item DC: 48.0
  \item UF: 770.0
  \item UC: 1816.5
\end{itemize}

DF needs a dramatically smaller search tree
than the other models. DC is clearly second-best, while UF and UC are much
weaker, and UC requires the largest tree by a wide margin.

This is consistent with the root-gap results:

\begin{itemize}
  \item better root relaxation $\Rightarrow$ fewer nodes,
  \item weaker root relaxation $\Rightarrow$ larger tree.
\end{itemize}

Compared with the
median values of DC, UF, and UC, DF uses dramatically fewer nodes. Therefore the good performance
of DF is not just due to a fast implementation: it is solving a much easier branch-and-bound problem.


\begin{figure}[h!]
\centering
\includegraphics[width=0.8\textwidth]{imgs_final/bb_nodes_boxplot.png}
\caption{Branch-and-bound nodes by formulation.}
\end{figure}

\subsection{Time spent for separation problems}

Separation time matters only for the cut formulations, because the flow
formulations do not solve separation subproblems in callbacks.

The median separation times are:

\begin{itemize}
  \item DC: \SI{0.725}{s}
  \item UC: \SI{18.467}{s}
\end{itemize}

The median share of total runtime spent in separation is:

\begin{itemize}
  \item DC: 0.242
  \item UC: 0.218
\end{itemize}

This means that UC spends much more absolute time in
separation than DC, but the \emph{fraction} of total runtime spent in separation
is almost the same between the two cut models. Therefore the weak
performance of UC is not explained only by ``callbacks are expensive''.
The more important problem is that UC has a much weaker root bound and a much
larger branch-and-bound tree. In other words:

\begin{quote}
Separation is not the main reason why UC is bad; the formulation itself is weak.
\end{quote}

\subsection{Dependence of the outcomes on the instance}

Table~\ref{tab:instance-dependence} summarizes the strongest observed
correlations between instance features and performance indicators. Overall, the
main drivers are the number of edges, graph density / average degree, and the
number of terminals.

\begin{table}[h!]
\centering
\small
\setlength{\tabcolsep}{5pt}
\begin{tabular}{llll}
    \toprule
    Formulation & Quantity & Main feature & Correlation \\
    \midrule
    DC & runtime & $n\_edges$ / density / avg.\ degree & $0.864$ \\
    DC & root LP gap & $n\_terminals$ & $-0.769$ \\
    DF & runtime & $n\_edges$ / density / avg.\ degree & $0.613$ \\
    DF & root LP gap & $n\_terminals$ & $-0.575$ \\
    UF & runtime & $n\_terminals$ & $0.646$ \\
    UF & root LP gap & $n\_edges$ / density / avg.\ degree & $0.637$ \\
    UC & separation time & $n\_terminals$ & $0.885$ \\
    UC & root LP gap & $n\_edges$ / density / avg.\ degree & $0.637$ \\
    \bottomrule
\end{tabular}
\caption{Strongest correlations between instance features and performance indicators.}
\label{tab:instance-dependence}
\end{table}

The table shows that larger and denser graphs are especially difficult for the
two directed formulations in terms of runtime. For DC the effect is very strong
(correlation $0.864$), while for the corrected DF it is still clearly visible
(correlation $0.613$). In contrast, UF is still most sensitive to the number of
terminals, which is consistent with its multi-commodity structure.

For the undirected cut formulation, the strongest
effect is the dependence of separation time on the number of terminals
(correlation $0.885$), indicating that separation becomes much more expensive on
instances with many terminals.

Another interesting observation is that for the
directed cut formulation the root LP gap is negatively correlated with the
number of terminals, suggesting that on this instance family its relaxation does
not deteriorate as the terminal set grows.

For DF, the reported root-gap correlations should again be interpreted
cautiously, because root-gap values are available only for a subset of the
instances. Still, the available data suggest that the corrected DF formulation
does not suffer from the same runtime bottlenecks as UF.


\subsection{More}

\subsubsection*{Median runtime alone is not enough}

Because there is a hard time limit of 120 seconds, median runtime by itself can
be misleading for the weaker formulations. In the updated benchmark this warning
mainly applies to UF and UC, whose medians are still essentially the time
limit. For DF and DC the medians are genuinely informative, because both
formulations solve all instances. Even so, the number of optimal solves and the
pairwise win counts remain more informative than the median runtime alone.


\section{Conclusion}

Across all five requested aspects, the same message appears repeatedly:

\begin{itemize}
  \item DF has the best overall runtime,
  \item DF has the strongest reported root LP behavior, although this should be
  interpreted cautiously because the callback does not always record a root
  bound before the solve ends,
  \item DF uses the fewest branch-and-bound nodes,
  \item DC keeps separation time under control while still benefiting from it,
  \item and the corrected DF implementation scales more favorably with instance difficulty than
  the other formulations.
\end{itemize}

Therefore the computational study strongly supports the directed flow
formulation as the best of the four on this benchmark set.
